% Flavor catalog per-process entry.
% process_id: K020
% family: kaon
% owner_agent_id: PKA-K020
% writer_agent_id: null
% checker_agent_id: null
% opus_signoff_id: null
% Canonical metadata sidecar: K020.yaml
% Status is controlled by the YAML sidecar status_history, not this comment.

\section{K020: \(K^+ \to \pi^+ e^{\pm}\mu^{\mp}\)}

\subsection*{Process}
\textbf{Standard notation:} \(K^+ \to \pi^+ e^{\pm}\mu^{\mp}\).  This entry
covers the charged-kaon lepton-flavor-violating semileptonic modes
\(K^+\to\pi^+\mu^+ e^-\) and \(K^+\to\pi^+\mu^- e^+\), process \(K020\) in
the kaon family.  It does not include the distinct lepton-number-violating
\(K^+\to\pi^-\mu^+e^+\) mode.

\subsection*{PDG / equivalent value(s)}
The canonical current limits are the PDG Live / PDG API 2025 charged-kaon
listing values, accessed 2026-05-17:
\[
  {\cal B}(K^+\to\pi^+\mu^+e^-) < 1.3\times10^{-11}\quad(90\%\ {\rm CL}),
\]
and
\[
  {\cal B}(K^+\to\pi^+\mu^-e^+) < 6.6\times10^{-11}\quad(90\%\ {\rm CL}).
\]
The local PDG snapshot is
\texttt{flavor\_catalog/references/K020/pdg2025\_kplus\_lfv\_semileptonic\_api.txt}.
PDG identifier \(\texttt{S010.29}\) is the \(\mu^+e^-\) charge assignment and
is based on the Sher et al. BNL E865 result combined with earlier E865 and
E777 data.  PDG identifier \(\texttt{S010.25}\) is the \(\mu^-e^+\) charge
assignment and uses the NA62 2021 result.

\subsection*{Relevance to the RS / anarchic-flavor pipeline}
This is a \textbf{SECONDARY} Wave-8 entry; see
\texttt{flavor\_catalog/PRIORITY\_TIERS.md} for the tier policy.  K020 is a
charged-kaon LFV companion to the neutral-kaon LFV channel K019.  At
the quark level it probes \(\Delta S=1\) semileptonic LFV transitions
\(s\to d\,e\mu\), not the already implemented \(\Delta F=2\) meson-mixing
lane.  In a quark-only anarchic RS scan it is outside the minimal implemented
surface, but in lepton-bulk or seesaw-extended RS models nonuniversal fermion
localization can generate off-diagonal quark and lepton couplings to KK gauge
bosons, shifted \(Z\) bosons, or Higgs/scalar states.  Those effects match onto
vector, scalar, or tensor four-fermion operators of the form
\((\bar s\Gamma d)(\bar e\Gamma'\mu)\).  Dipole LFV constraints such as
\(\mu\to e\gamma\) are correlated through the same lepton-sector spurions but
do not by themselves implement the \(K\to\pi e\mu\) rate.

\subsection*{Post-2008 developments}
The BNL E865 Sher et al. 2005 paper remains the canonical source for
\(K^+\to\pi^+\mu^+e^-\), with a \(2.1\times10^{-11}\) E865-only limit and
\(1.3\times10^{-11}\) after combination with earlier E865 and E777 data.
The key post-2008 experimental update is NA62 2021, arXiv:2105.06759, which
sets \({\cal B}(K^+\to\pi^+\mu^-e^+)<6.6\times10^{-11}\) and improves the
previous limit by about one order of magnitude for that charge assignment.

For post-CFW theory context, Beneke, Moch, and Rohrwild 2015
(arXiv:1508.01705) provide a full 5D RS treatment of charged-lepton flavor
violation, including dipole, \(\mu\to3e\), conversion, and custodial-model
effects.  Angelescu, Faroughy, and Sumensari 2020 (arXiv:2002.05684) give a
modern EFT treatment of quark-lepton LFV operators and explicitly list
low-energy semileptonic modes such as
\(K\to\pi \ell_k^\pm\ell_l^\mp\).  These sources matter because K020 cannot be
obtained by reusing a pure kaon-mixing Wilson coefficient or the existing
\(\mu\to e\gamma\) dipole check.

\subsection*{Constraint validity and model dependence}
Experimentally, these are clean upper limits on forbidden SM processes, so the
observable entry is an upper-limit constraint rather than an averaged
nonzero measurement.  The mapping to RS parameters is model dependent.  A
usable constraint needs a convention for the lepton bulk profiles or seesaw
spurions, the quark-lepton operator basis, and the \(K\to\pi\) hadronic form
factors for vector/scalar/tensor structures.  The two charge assignments should
be kept separate because models with chiral lepton couplings can populate them
differently.  Existing \(\mu\to e\gamma\), \(\mu\to3e\), and conversion limits
can dominate in many lepton-sector choices, but they do not eliminate the need
for a semileptonic kaon rate when quark-flavor violation is present.

\subsection*{Code coverage in this repo}
\textbf{NO.}  A repository search over \texttt{quarkConstraints/}, \texttt{qcd/},
\texttt{flavorConstraints/}, \texttt{neutrinos/}, \texttt{yukawa/},
\texttt{warpConfig/}, \texttt{solvers/}, \texttt{scanParams/}, and
\texttt{tests/}.  The K020-specific grep found only generic LFV metadata and
the existing \(\mu\to e\gamma\) dipole module, e.g.
\texttt{flavorConstraints/muToEGamma.py:75} and
\texttt{scanParams/scan.py:524}.  Broader greps find the implemented
\(\Delta F=2\) surfaces such as \texttt{quarkConstraints/deltaf2.py:755} for
\(\Delta m_K\), and the modern policy system list at
\texttt{quarkConstraints/modern/phenomenology.py:23} contains only
\(\varepsilon_K\), \(K\), \(B_d\), \(B_s\), and \(D0\).  There is no
\(K^+\to\pi^+e\mu\) observable, likelihood, operator basis, or test.

\subsection*{Implementation difficulty}
\textbf{HIGH.}  A production implementation would need a new
\(\Delta S=1\) semileptonic LFV operator basis, RS matching for simultaneous
quark and lepton flavor violation, and a \(K\to\pi e\mu\) rate calculation with
the relevant hadronic form factors.  The existing \(\Delta F=2\) and
\(\mu\to e\gamma\) code paths cannot be adapted by a threshold change or a
simple constant replacement.

\subsection*{Key references}
Process-local source keys before bibliography consolidation:
\texttt{PDG2025\_Kplus\_LFV\_semileptonic},
\texttt{Sher2005\_Kplus\_piplus\_mup\_em},
\texttt{NA622021\_Kplus\_piplus\_mum\_ep},
\texttt{CFW2008\_RS\_flavor},
\texttt{BenekeMochRohrwild2015\_RS\_LFV}, and
\texttt{AngelescuFaroughySumensari2020\_semileptonic\_LFV}.
